\chapter{어려운 방향의 증명 경로}
\label{chap:routes}

목표는 정리~\ref{thm:hard}이다. 이 장은 회전 체계 관점에서 두 증명 경로를 비교하고 하나를 고른다.

\section{제약: 교과서의 기하 논증을 조합적으로 다시 써야 한다}

이 문서에서 ``평면''은 종수 0인 회전 체계가 있다는 뜻이다 (정의~\ref{def:IsPlanar}). 교과서 증명은 다음과
같은 위상적 사실을 자유롭게 쓰지만, 우리는 그 하나하나를 면 순열의 언어로 다시 증명해야 한다.
\begin{itemize}
  \item 사이클은 평면을 안과 밖으로 나눈다 (조르당 곡선 정리).
  \item 2-연결 평면 그래프의 면 경계는 사이클이다.
  \item 한 면 안에 새 꼭짓점이나 변을 그려 넣을 수 있다.
\end{itemize}

\subsection*{지금 가진 도구로 할 수 있는 것}

쉬운 방향에서 만든 두 도구를 합치면 ``종수 0의 경직성''을 얻는다.
\begin{enumerate}
  \item 호환 하나를 곱하면 순환 수가 정확히 \(\pm 1\) 바뀐다 (보조정리~\ref{lem:swap}).
  \item 모든 회전 체계에서 \(V - E + F \le 2c\)이다 (정리~\ref{thm:euler}).
\end{enumerate}
\textbf{결론}: 연결 성분 수는 그대로 두면서 \(V - E + F\)를 2 늘리는 수술은 종수 0 회전 체계에서는
불가능하다. 그런 수술이 가능해 보이는 상황은 모순이다.

예를 들어 2-연결 평면 그래프의 한 면이 꼭짓점 \(w\)를 두 번 지난다고 하자. 그 두 모서리에서 \(w\)를 둘로
쪼개면 꼭짓점이 1, 면이 1 늘고 성분은 그대로다. 이것이 위의 결론과 모순이므로 ``면 경계는 사이클''이 따라
나온다 (보조정리~\ref{lem:J1}). 이처럼 \emph{국소적인} 기하 사실은 이 방법으로 얻을 수 있다. 반면 ``이
사이클의 안쪽에 있는 브리지들'' 같은 \emph{전역적인} 구조는 이 방법으로 다루기 어렵다. 경로를 고를 때
가장 중요한 기준이 이것이다.

이 형식화는 Mathlib의 단순 그래프 위에 지지 집합 기준 꼭짓점 연결도(\(k\)-연결)와
꼭짓점 타입을 유지하는 변 수축을 정의한다. 주 증명은 그래프 마이너 대신 세분을 직접 다룬다.

\section{경로 A: 변 수축 (Thomassen 1981)}

C.~Thomassen, \emph{Kuratowski's theorem}, J.~Graph Theory 5 (1981). 같은 구조가 Diestel,
\emph{Graph Theory}, 4.4절에도 있다 (Diestel은 마이너를 쓴다).
\begin{enumerate}
  \item 최소 반례는 3-연결이다. 크기 2 이하의 분리가 있으면 양쪽을 따로 평면으로 만들어 붙인다
        (정리~\ref{thm:reduce}).
  \item 꼭짓점이 5개 이상인 3-연결 그래프에는, 수축해도 3-연결인 변 \(xy\)가 있다
        (정리~\ref{thm:contractible}).
  \item \(G\)에 쿠라토프스키 부분 그래프가 없으면 \(G/xy\)에도 없다 (보조정리~\ref{lem:lift}).
  \item 귀납법으로 \(G/xy\)는 평면이다. 합친 꼭짓점 \(v\)를 지운 \(G/xy - v\)에서, 원래 \(v\)가 있던 면은
        \(v\)의 이웃을 모두 지나는 사이클 \(C\)로 둘러싸인다 (보조정리~\ref{lem:J2}).
  \item \(C\)를 따라 \(x\)의 이웃과 \(y\)의 이웃이 어떻게 놓였는지 본다 (보조정리~\ref{lem:cyclic}).
    \begin{itemize}
      \item 서로 섞이지 않으면 \(v\)를 \(x\), \(y\)로 다시 나눠 \(G\)의 평면 회전 체계를 만든다
            (보조정리~\ref{lem:S}).
      \item 공통 이웃이 3개 이상이면 \(G\) 안에 \(TK_5\)가 있다.
      \item 번갈아 놓이면 \(G\) 안에 \(TK_{3,3}\)이 있다.
    \end{itemize}
\end{enumerate}

\subsection*{회전 체계로 옮길 때}
\begin{itemize}
  \item 4는 ``꼭짓점을 지우면 그 둘레의 면들이 한 면으로 합쳐진다''는 계산과 보조정리~\ref{lem:J1}로
        해결된다. 모두 꼭짓점 \(v\) 근처의 국소적인 계산이다.
  \item 5의 ``나누기''는 \(v\)의 회전 순서를 두 구간으로 자르는 국소적인 수술이다. 면 수는 호환
        보조정리로 센다.
  \item 5의 경우 나누기는 순환 수열에 대한 순수한 조합 문제다.
  \item 기하가 남는 곳은 4뿐이고 그것도 국소적이다.
\end{itemize}
\textbf{새로 세울 이론}: 3-연결성, 변 수축, 0/1/2-합 붙이기, 세분의 경로 표현, 들어 올리기 보조정리.

\section{경로 B: 사이클과 브리지}

Bondy--Murty식 증명, 또는 평면성 판정 알고리즘의 구조다.
\begin{enumerate}
  \item 최소 비평면 그래프 \(G\)에서 변 \(e = uv\)를 하나 빼면 평면이다.
  \item \(u\), \(v\)를 지나는 사이클 가운데 ``안쪽''이 가장 큰 \(C\)를 고른다.
  \item \(C\)의 브리지를 안쪽과 바깥쪽으로 나누고, 서로 겹치는(overlap) 브리지의 배치를 분석해 \(K_5\)나
        \(K_{3,3}\)을 찾는다.
\end{enumerate}

\subsection*{회전 체계로 옮길 때}
\begin{itemize}
  \item 임의의 사이클 \(C\)에 대해 면과 브리지를 ``안쪽/바깥쪽''으로 나누는 조합적 조르당 곡선 정리가
        필요하다. 이것은 전역적인 진술이다.
  \item ``안쪽이 가장 크다''는 극값 선택은 안쪽의 크기를 정의하고 비교해야 한다.
  \item 브리지 겹침 이론과 마지막 경우 나누기(대략 5가지)가 경로 A보다 무겁다.
\end{itemize}

\section{비교와 결정}

\begin{center}
\begin{tabular}{lll}
\hline
 & 경로 A (수축) & 경로 B (브리지) \\
\hline
필요한 기하 사실 & 면 경계가 사이클이다 (국소적) & 사이클의 안팎 분리 (전역적) \\
기하를 옮기는 방법 & 종수 0의 경직성 (이미 가진 도구) & 새 이론이 필요 \\
새로 세울 그래프 이론 & 3-연결성, 수축, 붙이기 & 브리지, 겹침, 극값 사이클 \\
마지막 경우 나누기 & 순환 수열 3가지 (전수 검사 완료) & 브리지 배치 약 5가지 \\
지금 도구 재사용 & 변 지우기, 세분, 호환 보조정리 & 변 지우기 \\
\hline
\end{tabular}
\end{center}

\textbf{경로 A를 고른다.} 세부 선택은 다음과 같다.
\begin{itemize}
  \item \textbf{마이너 대신 위상 소수.} Diestel처럼 마이너를 쓰면 마이너를 정의하고 ``\(K_5\) 마이너면
        \(TK_5\) 또는 \(TK_{3,3}\)''을 따로 증명해야 한다. Thomassen의 원래 증명처럼 처음부터 세분(위상
        소수)만 쓰고, 수축에 대해서는 들어 올리기 보조정리 하나로 처리한다.
  \item \textbf{세분은 경로 표현으로 다룬다.} 들어 올리기와 \(TK_5\)/\(TK_{3,3}\) 구성은 ``분기점 +
        내부가 서로소인 경로들''로 생각해야 자연스럽다. 기존 \texttt{IsSubdivision}(귀납 정의)을 유지하고,
        경로 표현과 동치임을 한 번 증명한다 (정리~\ref{thm:paths_iff}).
\end{itemize}
구현한 설계 결정은 \ref{chap:plan}장에 있다.

\section{규모 감각}

쉬운 방향(Lean 약 2,200줄)에 비해 2--4배로 어림한다. 근거는 새로 세울 이론의 수(3-연결성, 수축, 붙이기,
경로 표현)이고, 정밀한 추정은 아니다. 가장 큰 불확실성은 2-합 붙이기와 꼭짓점 나누기의 다트 대응
작업이다. 세분 되돌리기(정리~\ref{thm:smooth})에서 비슷한 작업이 약 360줄이었다.
